\section*{Hoofdstuk \ref{ch:g7:prop} --- Evenredigheid}

\begin{solution}{exo:g7:prop:1}
Eerste tabel: \omterm{def:g3:division:remainder}{quotiënten} $\frac{7.5}{3} = 2.5$, $\frac{17.5}{7} = 2.5$,
$\frac{27.5}{11} = 2.5$: \omterm{def:g6:propdata:def}{evenredig}, met constante $2.5$.

Tweede tabel: $\frac62 = 3$ en $\frac{15}{5} = 3$, maar
$\frac{28}{9} \approx 3.1$: niet \omterm{def:g6:propdata:def}{evenredig}.
\end{solution}

\begin{solution}{exo:g7:prop:2}
Met de kruisregel: $5 \times p = 6 \times 8$, dus $p = \frac{48}{5} = 9.60$ euro.
\end{solution}

\begin{solution}{exo:g7:prop:3}
Tempo: $36 \div 3 = 12$ bladzijden per minuut. In $5$ minuten: $60$ bladzijden.
Voor $96$ bladzijden: $96 \div 12 = 8$ minuten.
\end{solution}

\begin{solution}{exo:g7:prop:4}
$30\,\%$ van $250$: $0.3 \times 250 = 75$. \quad
$15\,\%$ van $60$: $0.15 \times 60 = 9$. \quad
$t\,\%$ van $80$ is $20$: $\frac{20}{80} = \frac{25}{100}$, dus $t = 25$.
\end{solution}

\begin{solution}{exo:g7:prop:5}
$\frac{18}{25} = \frac{18 \times 4}{25 \times 4} = \frac{72}{100} = 72\,\%$.
\end{solution}

\begin{solution}{exo:g7:prop:6}
Pad: $9 \times 25\,000 = 225\,000$ cm $= 2.25$ km. Meer: $2$ km $= 200\,000$ cm op
de grond, dus $200\,000 \div 25\,000 = 8$ cm op de kaart.
\end{solution}

\begin{solution}{exo:g7:prop:7}
$4.3$ m $= 430$ cm, en $430 \div 43 = 10$ cm.
\end{solution}

\begin{solution}{exo:g7:prop:8}
Chocolade per portie: $240 \div 8 = 30$ g. Met $300$ g: $300 \div 30 = 10$ porties
precies.
\end{solution}

\begin{solution}{exo:g7:prop:9}
\emph{1.} Korting: $20$ euro, dus $\frac{20}{80} = 25\,\%$ van de oorspronkelijke
prijs.

\emph{2.} De stijging van $20$ euro wordt nu vergeleken met het \emph{nieuwe}
uitgangspunt $60$: $\frac{20}{60} = \frac13 \approx 33\,\%$. Een percentage
verwijst altijd naar een geheel; dat geheel is veranderd, en dus het percentage
ook.
\end{solution}

\begin{solution}{exo:g7:prop:10}
De constante uit de volledige kolom: $\frac{17.5}{10} = 1.75$. Dan is
$y = 4 \times 1.75 = 7$ en $x = \frac{10.5}{1.75} = 6$.
\end{solution}

\begin{solution}{exo:g7:prop:11}
\emph{1.} Na één kopie: $10 \times 0.8 = 8$ cm. Na twee: $8 \times 0.8 = 6.4$ cm.

\emph{2.} De tweede verkleining werkt op de al verkleinde kopie, niet op het
origineel: de factoren vermenigvuldigen zich, $0.8 \times 0.8 = 0.64$, dus
verkleinen twee kopieën tot $64\,\%$ van het origineel --- geen $60\,\%$.
\end{solution}

\begin{solution}{pb:g7:prop:1}
\textbf{1.} Hoogtes en schaduwen zijn \omterm{def:g6:propdata:def}{evenredig}, dus geldt volgens de kruisregel
(\cref{thm:g7:prop:cross}), met de kolom van de stok $(1,\ 0.4)$ en die van de
boom $(h,\ 3.2)$: $0.4 \times h = 1 \times 3.2$, dus
$h = \frac{3.2}{0.4} = 8$ m.

\textbf{2.} Toren: $\frac{26}{0.4} = 65$ m. Schaduw van het kind:
$1.5 \times 0.4 = 0.6$ m (schaduw $=$ hoogte $\times$ de constante $0.4$ van dat
moment).

\textbf{3.} Terwijl de zon langs de hemel trekt, verandert de constante die hoogtes
met schaduwen verbindt; alleen metingen op hetzelfde moment delen dezelfde
constante.

\textbf{4.} Op dat moment is de constante $1$: elke hoogte is gelijk aan haar
schaduw. De top van de \omterm{def:g5:solids:def}{piramide} is dus $147$ m hoog --- de hoogte van de Grote
\omterm{def:g5:solids:def}{Piramide} (de punt van haar schaduw, gemeten vanaf het punt onder de top, is alles
wat je nodig hebt).

\textbf{5.} Het kolomquotiënt $\frac{\text{schaduw}}{\text{hoogte}}$ --- de
evenredigheidsconstante van dat moment --- is voor elk verticaal voorwerp gelijk:
en dat is precies wat de tabel van hoogtes en schaduwen een \omterm{def:g7:prop:table}{evenredigheidstabel}
maakt (\cref{def:g7:prop:table}).

\textbf{6.} Zonnestralen komen \omterm{def:g6:lines:perp}{parallel} aan. Op een \emph{platte} aarde zouden twee
parallelle stralen twee verticale stokken onder dezelfde hoek raken: beide zouden
evenredige schaduwen werpen --- of beide geen schaduw, of beide wel. Eén stok
zonder schaduw (Syene) en één met schaduw (Alexandrië), op hetzelfde ogenblik, is
op een platte aarde onmogelijk: de twee verticalen moeten in verschillende
richtingen wijzen --- het oppervlak kromt.

\textbf{7.} Op de tekening wijst de verticaal van Syene recht naar de zon; die van
Alexandrië staat schuin, over de hoek tussen de richtingen van de twee steden
gezien vanuit het middelpunt. Omdat de stralen \omterm{def:g6:lines:perp}{parallel} zijn, is die scheefstand
precies de $7.2^\circ$ die in Alexandrië tussen straal en stok gemeten is.

\textbf{8.} $\frac{7.2}{360} = \frac{72}{3600} = \frac{1}{50}$: een vijftigste van
een volle draai.

\textbf{9.} Booglengtes zijn \omterm{def:g6:propdata:def}{evenredig} met hoeken: hoort $7.2^\circ$ --- een
vijftigste van de draai --- bij $800$ km, dan hoort de hele draai bij
\[
50 \times 800 = 40\,000 \text{ km}.
\]

\textbf{10.} \omterm{def:g4:geometry:circle}{Diameter} $= \frac{\text{omtrek}}{\pi} \approx
\frac{40\,000}{3.14} \approx 12\,700$ km --- tegen de moderne $12\,742$ km: ruim
binnen één procent juist, met een stok.

\textbf{11.} $40\,000$ km $= 4\,000\,000\,000$ cm, en
$4\,000\,000\,000 \div 40\,000\,000 = 100$ cm $= 1$ m \omterm{def:g6:measure:perimeter}{omtrek}. \omterm{def:g4:geometry:circle}{Diameter}:
$100 \div 3.14 \approx 32$ cm --- een mooie bureauglobe.

\textbf{12.} $4.8$ km $= 480\,000$ cm, en
$480\,000 \div 40\,000\,000 = 0.012$ cm $= 0.12$ mm. De hoogste berg van de Alpen
is een tiende millimeter op een globe met een \omterm{def:g6:measure:perimeter}{omtrek} van een meter: een stofje. Op
deze \omterm{def:g7:prop:scale}{schaal} is de aarde gladder dan de meeste gepolijste ballen.

\textbf{13.} $40\,000 \div 40 = 1\,000$ dagen, en $1\,000 \div 365 \approx 2.7$:
bijna drie jaar lopen.

\textbf{14.} $\frac{10}{6\,370} \approx 0.0016$, dus ongeveer $0.2\,\%$
(preciezer $0.16\,\%$): de atmosfeer is een velletje, dun als een zucht, om de
planeet.

\textbf{15.} Fout: $42\,000 - 40\,000 = 2\,000$ km, en
$\frac{2\,000}{40\,000} = \frac{5}{100} = 5\,\%$. In één zin: met een stok, een
put, een opgemeten weg en één \omterm{def:g7:prop:table}{evenredigheid} mat Eratosthenes een planeet op een paar
procent nauwkeurig --- wiskunde komt heel ver met heel weinig.
\end{solution}
